\documentclass[10pt,a4paper]{amsart}
\usepackage{amsmath,amssymb,amsthm,graphicx}
% Retained minimal preamble for the Stacks Project chapter more-algebra.tex (More on Algebra).
% The chapter file pinned-original-source.tex carries only \input{preamble} on its first line, so no
% preamble text lives inside the chapter itself.  The Stacks Project uses one single shared preamble for
% every chapter; the snapshot stages it identically as sources/stacks-chow/preamble.tex and
% sources/stacks-etale-cohomology/preamble.tex.  The block below is copied mechanically, line for line,
% from sources/stacks-chow/preamble.tex lines 14-18 (the Xy-pic setup) and 154-217 (the native theorem
% declarations together with the mathematical macros).  Nothing was added, removed or redefined; the
% wrapper class and the amsmath/amssymb/amsthm/graphicx base packages are supplied by the assembler.
% Chapter: more-algebra.tex -- More on Algebra.
% For commutative diagrams we use Xy-pic
\usepackage[all]{xy}

% We use 2cell for 2-commutative diagrams.
\xyoption{2cell}

% Theorem environments.
%
\theoremstyle{plain}
\newtheorem{theorem}[subsection]{Theorem}
\newtheorem{proposition}[subsection]{Proposition}
\newtheorem{lemma}[subsection]{Lemma}

\theoremstyle{definition}
\newtheorem{definition}[subsection]{Definition}
\newtheorem{example}[subsection]{Example}
\newtheorem{exercise}[subsection]{Exercise}
\newtheorem{situation}[subsection]{Situation}

\theoremstyle{remark}
\newtheorem{remark}[subsection]{Remark}
\newtheorem{remarks}[subsection]{Remarks}

\numberwithin{equation}{subsection}

% Macros
%
\def\lim{\mathop{\mathrm{lim}}\nolimits}
\def\colim{\mathop{\mathrm{colim}}\nolimits}
\def\Spec{\mathop{\mathrm{Spec}}}
\def\Hom{\mathop{\mathrm{Hom}}\nolimits}
\def\Ext{\mathop{\mathrm{Ext}}\nolimits}
\def\SheafHom{\mathop{\mathcal{H}\!\mathit{om}}\nolimits}
\def\SheafExt{\mathop{\mathcal{E}\!\mathit{xt}}\nolimits}
\def\Sch{\mathit{Sch}}
\def\Mor{\mathop{\mathrm{Mor}}\nolimits}
\def\Ob{\mathop{\mathrm{Ob}}\nolimits}
\def\Sh{\mathop{\mathit{Sh}}\nolimits}
\def\NL{\mathop{N\!L}\nolimits}
\def\CH{\mathop{\mathrm{CH}}\nolimits}
\def\proetale{{pro\text{-}\acute{e}tale}}
\def\etale{{\acute{e}tale}}
\def\QCoh{\mathit{QCoh}}
\def\Ker{\mathop{\mathrm{Ker}}}
\def\Im{\mathop{\mathrm{Im}}}
\def\Coker{\mathop{\mathrm{Coker}}}
\def\Coim{\mathop{\mathrm{Coim}}}

% Boxtimes
%
\DeclareMathSymbol{\boxtimes}{\mathbin}{AMSa}{"02}

%
% Macros for moduli stacks/spaces
%
\def\QCohstack{\mathcal{QC}\!\mathit{oh}}
\def\Cohstack{\mathcal{C}\!\mathit{oh}}
\def\Spacesstack{\mathcal{S}\!\mathit{paces}}
\def\Quotfunctor{\mathrm{Quot}}
\def\Hilbfunctor{\mathrm{Hilb}}
\def\Curvesstack{\mathcal{C}\!\mathit{urves}}
\def\Polarizedstack{\mathcal{P}\!\mathit{olarized}}
\def\Complexesstack{\mathcal{C}\!\mathit{omplexes}}
% \Pic is the operator that assigns to X its picard group, usage \Pic(X)
% \Picardstack_{X/B} denotes the Picard stack of X over B
% \Picardfunctor_{X/B} denotes the Picard functor of X over B
\def\Pic{\mathop{\mathrm{Pic}}\nolimits}
\def\Picardstack{\mathcal{P}\!\mathit{ic}}
\def\Picardfunctor{\mathrm{Pic}}
\def\Deformationcategory{\mathcal{D}\!\mathit{ef}}
\usepackage{cleveref}
\crefname{originallemma}{original lemma}{original lemmas}
\Crefname{originallemma}{Original lemma}{Original lemmas}
\makeatletter
\newlabel{algebra-lemma-adjoint-hom-restrict}{{08YP}{source}{}{sourcepointer-1}{}}
\newlabel{algebra-lemma-adjoint-hom-restrict@cref}{{[originallemma][0][]08YP}{[1][0][]source}{}{}{}}
\makeatother
\crefname{originallemma}{original lemma}{original lemmas}
\Crefname{originallemma}{Original lemma}{Original lemmas}
\makeatletter
\newlabel{derived-lemma-K-injective}{{070I}{source}{}{sourcepointer-2}{}}
\newlabel{derived-lemma-K-injective@cref}{{[originallemma][0][]070I}{[1][0][]source}{}{}{}}
\makeatother
\crefname{originallemma}{original lemma}{original lemmas}
\Crefname{originallemma}{Original lemma}{Original lemmas}
\makeatletter
\newlabel{derived-section-derived-limit}{{08TB}{source}{}{sourcepointer-3}{}}
\newlabel{derived-section-derived-limit@cref}{{[originallemma][0][]08TB}{[1][0][]source}{}{}{}}
\makeatother
\crefname{originalequation}{original equation}{original equations}
\Crefname{originalequation}{Original equation}{Original equations}
\makeatletter
\newlabel{equation-first-ss-ext}{{17370}{source}{}{sourcepointer-4}{}}
\newlabel{equation-first-ss-ext@cref}{{[originalequation][0][]17370}{[1][0][]source}{}{}{}}
\makeatother
\crefname{originallemma}{original lemma}{original lemmas}
\Crefname{originallemma}{Original lemma}{Original lemmas}
\makeatletter
\newlabel{homology-section-double-complex}{{012X}{source}{}{sourcepointer-5}{}}
\newlabel{homology-section-double-complex@cref}{{[originallemma][0][]012X}{[1][0][]source}{}{}{}}
\makeatother
\crefname{originallemma}{original lemma}{original lemmas}
\Crefname{originallemma}{Original lemma}{Original lemmas}
\makeatletter
\newlabel{lemma-hom-K-injective}{{13956}{source}{}{sourcepointer-6}{}}
\newlabel{lemma-hom-K-injective@cref}{{[originallemma][0][]13956}{[1][0][]source}{}{}{}}
\makeatother
\title{Eight equivalent conditions for the vanishing of Ext groups out of a localization at a single element}
\author{The Stacks Project Authors}
\date{stacks-project:more-algebra.tex:local-snapshot-20261012; GNU FDL 1.2 or later}
\begin{document}
\maketitle
\setcounter{subsection}{1}
\noindent\textit{Editorial note.} All delivered mathematical text is the author's own, copied from the pinned Stacks Project chapter file more-algebra.tex; no statement or proof step was rewritten, reordered or removed.\par\medskip
\noindent\textit{Editorial note.} The chapter file pinned-original-source.tex carries only \textbackslash{}input\{preamble\} on its first line, so no preamble text lives inside the chapter. The retained file sources/stacks-more-algebra/retained-preamble.tex is a mechanical, line-for-line copy of the project's single shared preamble as staged in the snapshot; nothing was added, removed or redefined and no package was added.\par\medskip
\noindent\textit{Editorial note.} The source URL is the source-level Stacks Project URL, as required when more than one card is produced from one source; the per-claim tag is carried by claim\_key.\par\medskip
\noindent\textit{Editorial note.} The last enumerated item of the statement, 'add more here.', is the author's own placeholder in the pinned source; it is delivered verbatim rather than removed, and it contributes no mathematical condition. The footnote in the counted proof span, which explains the unbounded-below case and ends 'Details omitted.', is likewise delivered verbatim; that footnote records the author applying the argument to a replacement double complex rather than invoking the referenced equation, and it is part of the counted proof span. The author's typo 'we chooses' is reproduced unchanged.\par\medskip
\noindent\textit{Editorial note.} The references inside the counted proof span are recorded as non-load-bearing navigation pointers: algebra-lemma-adjoint-hom-restrict, derived-lemma-K-injective, derived-section-derived-limit and homology-section-double-complex resolve to the project's cross-chapter tag list, while lemma-hom-K-injective and equation-first-ss-ext are same-chapter labels.\par\medskip
\noindent\textit{Allowed standard prerequisites (editorial).}\begin{itemize}
\item K-injective complexes and the fact that an acyclic K-injective complex is homotopy equivalent to zero. used at proof lines 25870-25884. Is the device that turns the vanishing of Ext groups into the vanishing of maps in the derived category.
\item The adjunction identifying Hom over A from A\_f with Hom over A\_f after applying Hom\_A(A\_f, -). used at proof lines 25875-25879. Lets a map out of an A\_f-complex be rewritten as a map into Hom\_A(A\_f, I\textasciicircum{}bullet).
\item The explicit two-term free resolution of A\_f and the construction of derived limits. written out at proof lines 25887-25895 and invoked at proof lines 25902-25905. Supplies the displayed pro-product map in condition (7) and the identification of T(K, f) with R Hom\_A(A\_f, K).
\item Ext spectral sequence of a two-term projective resolution and its degeneration. used at proof lines 25909-25925. Produces the short exact sequences that relate Ext\textasciicircum{}p\_A(A\_f, K) to Hom and Ext\textasciicircum{}1 of the cohomology modules.
\end{itemize}\medskip
\ifdefined\hypertarget\hypertarget{sourcepointer-1}{}\fi
\noindent\textit{Source reference.} Original lemma 08YP: focused-seeds/AGP/sources/stacks-snapshot-20261011/tags.txt, source lines 11347--11347. Cross-chapter navigation pointer into the Stacks Project; the standard fact it supplies is applied inside the counted proof span and the argument resting on it stays there.\par\medskip
\ifdefined\hypertarget\hypertarget{sourcepointer-2}{}\fi
\noindent\textit{Source reference.} Original lemma 070I: focused-seeds/AGP/sources/stacks-snapshot-20261011/tags.txt, source lines 8957--8957. Cross-chapter navigation pointer into the Stacks Project; the standard fact it supplies is applied inside the counted proof span and the argument resting on it stays there.\par\medskip
\ifdefined\hypertarget\hypertarget{sourcepointer-3}{}\fi
\noindent\textit{Source reference.} Original lemma 08TB: focused-seeds/AGP/sources/stacks-snapshot-20261011/tags.txt, source lines 11159--11159. Cross-chapter navigation pointer into the Stacks Project; the standard fact it supplies is applied inside the counted proof span and the argument resting on it stays there.\par\medskip
\ifdefined\hypertarget\hypertarget{sourcepointer-4}{}\fi
\noindent\textit{Source reference.} Original equation 17370: sources/stacks-more-algebra/pinned-original-source.tex, source lines 17370--17370. Same-chapter navigation pointer; the statement it names is standard for the section and the argument that consumes it stays inside the counted proof span.\par\medskip
\ifdefined\hypertarget\hypertarget{sourcepointer-5}{}\fi
\noindent\textit{Source reference.} Original lemma 012X: focused-seeds/AGP/sources/stacks-snapshot-20261011/tags.txt, source lines 1415--1415. Cross-chapter navigation pointer into the Stacks Project; the standard fact it supplies is applied inside the counted proof span and the argument resting on it stays there.\par\medskip
\ifdefined\hypertarget\hypertarget{sourcepointer-6}{}\fi
\noindent\textit{Source reference.} Original lemma 13956: sources/stacks-more-algebra/pinned-original-source.tex, source lines 13956--13956. Same-chapter navigation pointer; the statement it names is standard for the section and the argument that consumes it stays inside the counted proof span.\par\medskip
\begin{lemma}
\label{lemma-hom-from-Af}
Let $A$ be a ring. Let $f \in A$. Let $K \in D(A)$.
The following are equivalent
\begin{enumerate}
\item $\Ext^n_A(A_f, K) = 0$ for all $n$,
\item $\Hom_{D(A)}(E, K) = 0$ for all $E$ in $D(A_f)$,
\item $T(K, f) = 0$,
\item for every $p \in \mathbf{Z}$ we have $T(H^p(K), f) = 0$,
\item for every $p \in \mathbf{Z}$ we have
$\Hom_A(A_f, H^p(K)) = 0$ and $\Ext^1_A(A_f, H^p(K)) = 0$,
\item $R\Hom_A(A_f, K) = 0$,
\item the map $\prod_{n \geq 0} K \to \prod_{n \geq 0} K$,
$(x_0, x_1, \ldots) \mapsto (x_0 - fx_1, x_1 - fx_2, \ldots)$
is an isomorphism in $D(A)$, and
\item add more here.
\end{enumerate}
\end{lemma}

\begin{proof}
It is clear that (2) implies (1) and that (1) is equivalent to (6). Assume (1).
Let $I^\bullet$ be a K-injective complex of $A$-modules representing $K$.
Condition (1) signifies that $\Hom_A(A_f, I^\bullet)$ is acyclic.
Let $M^\bullet$ be a complex of $A_f$-modules representing $E$.
Then
$$
\Hom_{D(A)}(E, K) =
\Hom_{K(A)}(M^\bullet, I^\bullet) =
\Hom_{K(A_f)}(M^\bullet, \Hom_A(A_f, I^\bullet))
$$
by Algebra, Lemma \ref{algebra-lemma-adjoint-hom-restrict}.
As $\Hom_A(A_f, I^\bullet)$ is a K-injective complex of
$A_f$-modules by Lemma \ref{lemma-hom-K-injective}
the fact that it is acyclic implies that it is homotopy equivalent to zero
(Derived Categories, Lemma \ref{derived-lemma-K-injective}).
Thus we get (2).

\medskip\noindent
A free resolution of the $A$-module $A_f$ is given by
$$
0 \to \bigoplus\nolimits_{n \in \mathbf{N}} A \to
\bigoplus\nolimits_{n \in \mathbf{N}} A
\to A_f \to 0
$$
where the first map sends the $(a_0, a_1, a_2, \ldots)$ to
$(a_0, a_1 - fa_0, a_2 - fa_1, \ldots)$ and the second map sends
$(a_0, a_1, a_2, \ldots)$ to $a_0 + a_1/f + a_2/f^2 + \ldots$.
Applying $\Hom_A(-, I^\bullet)$ we get
$$
0 \to \Hom_A(A_f, I^\bullet) \to \prod I^\bullet \to \prod I^\bullet \to 0
$$
Since $\prod I^\bullet$ represents $\prod_{n \geq 0} K$
this proves the equivalence of (1) and (7). On the other hand,
by construction of derived limits in
Derived Categories, Section \ref{derived-section-derived-limit}
the displayed exact sequence shows the object $T(K, f)$
is a representative of $R\Hom_A(A_f, K)$ in $D(A)$.
Thus the equivalence of (1) and (3).

\medskip\noindent
There is a spectral sequence
$$
E_2^{p, q} = \Ext^p_A(A_f, H^q(K)) \Rightarrow
\Ext^{p + q}_A(A_f, K)
$$
See Equation (\ref{equation-first-ss-ext})\footnote{If $K$ is not
bounded below, we can't use the given reference.
In that case, we chooses a complex $K^\bullet$ representing $K$
and a free two-term resolution $0 \to F^{-1} \to F^0 \to A_f \to 0$
(see above). Then we get a double complex with terms $\Hom_A(F^{-q}, K^p)$
whose totalization computes $R\Hom(A_f, K)$. The second spectral sequence
associated to this double complex (Homology, Section
\ref{homology-section-double-complex})
gives us what we want. Details omitted.}.
This spectral sequence degenerates at $E_2$ because
$A_f$ has a length $1$ resolution by projective $A$-modules (see above)
hence the $E_2$-page has only 2 nonzero columns. Thus we obtain short exact
sequences
$$
0 \to \Ext^1_A(A_f, H^{p - 1}(K)) \to
\Ext^p_A(A_f, K) \to
\Hom_A(A_f, H^p(K)) \to 0
$$
This proves (4) and (5) are equivalent to (1).
\end{proof}

\end{document}
